%% ===== Figure_1.png  label=fig:arch
Reference architecture. Decisions are made per service.
ETSI-014-inspired paired-KME behaviour is isolated from protocol-specific
adapters. MACsec and OTN adapters are generic analytical abstractions, not
claims of standardised PQC/QKD bindings. Bidirectional arrows denote
policy/key requests downward and telemetry/status upward. \rev{No carrier data
are used anywhere in the evaluation.}

%% ===== Figure_2.png  label=fig:qkd
Attenuation-envelope rate and route qualification. (a) Single-span
envelope with the qualified direct spans and the per-hop operating points of
the two qualified trusted-node core routes. (b) Qualification outcome of the
six long-haul routes; rejected routes are listed with the constraint that
rejected them so that qualification failure is not mistaken for a measured
secret-key rate.

%% ===== Figure_3.png  label=fig:mix
Mode assignments per service rather than per physical link. QKD
appears only after route qualification and only for high-assurance services.

%% ===== Figure_4.png  label=fig:pool
(a) Pool trajectories for the two outage-affected routes and one
unaffected DCI route; shading is link-specific. (b) \rev{The predicted
reserve-time signal of Eq.~\eqref{eq:reserve} with the gold demotion margin and
the re-entry threshold.} The DCI outage triggers
reserve-time demotion; the shorter core outage does not.

%% ===== Figure_5.png  label=fig:latency
Per-session sampled rekey-control latency by protocol adapter
\rev{with the configured latency budgets (dotted)}. The
curves are analytical and seeded; they are not measured carrier p99 values.

%% ===== Figure_6.png  label=fig:tco
Illustrative cost model with physical-hop accounting. (a) \rev{Shared
versus incremental five-year NPV (Table~\ref{tab:costsplit}).} (b) Capex
sensitivity: common orchestration is charged to both strategies and QKD pairs
are counted per physical hop; the shaded area is the incremental assurance
premium. (c) \rev{Premium versus discount rate for three horizons.}

%% ===== Figure_7.png  label=fig:oat
\rev{One-at-a-time sensitivity of the physical QKD footprint (top), the
QKD-backed share of QKD-preferred rekeys (middle), and the minimum pool level
during the 24-hour run (bottom) to the six route-qualification and reserve
parameters. The reference value is outlined. Cliffs in the reserve parameters
correspond to violations of the dimensioning condition, Eq.~\eqref{eq:dimensioning}.}

%% ===== Figure_8.png  label=fig:grid
\rev{Joint sensitivity to the designed hop limit and the trusted-node
policy limit: (a) physical QKD hops with the number of overlaid logical links
in brackets; (b) five-year assurance premium. The reference configuration is
outlined.}

%% ===== Figure_9.png  label=fig:baselines
\rev{Comparison of the proposed policy with five alternative planning
and gating schemes on the reference scenario (Table~\ref{tab:baselines}).}

%% ===== Figure_10.png  label=fig:seeds
\rev{Distribution of the headline results over 20 seeded topologies
with the reference seed highlighted. Boxes show quartiles and whiskers the
range.}
